% Part: sets-functions-relations
% Chapter: functions
% Section: partial-functions

\documentclass[../../../include/open-logic-section]{subfiles}


\begin{document}

\olfileid{sfr}{fun}{par}

\olsection{আংশিক অপেক্ষক}

\begin{explain}
কখনো কখনো অপেক্ষকের সংজ্ঞাটি শিথিল করা উপযোগী হয়, যাতে সকল সম্ভাব্য ইনপুটের জন্য অপেক্ষকের আউটপুট সংজ্ঞায়িত থাকা আবশ্যক না হয়। এই ধরনের চিত্রণকে \emph{আংশিক অপেক্ষক} বলা হয়।
\end{explain}

\begin{defn}
একটি \emph{আংশিক অপেক্ষক} $f \colon A \pto B$ হলো এমন একটি চিত্রণ, যা~$A$-এর প্রত্যেক !!{element}-এর জন্য~$B$-এর বড়জোর একটি !!{element} নির্দিষ্ট করে। যদি $f$, $x \in A$-এর জন্য~$B$-এর কোনো উপাদান নির্দিষ্ট করে, তবে বলি $f(x)$ \emph{সংজ্ঞায়িত}; অন্যথায় \emph{অসংজ্ঞায়িত}। $f(x)$ সংজ্ঞায়িত হলে আমরা $f(x) \fdefined$ লিখি, অন্যথায় $f(x) \fundefined$ লিখি। একটি আংশিক অপেক্ষক~$f$-এর \emph{সংজ্ঞাক্ষেত্র} হলো~$A$-এর সেই উপসেট, যেখানে সেটি সংজ্ঞায়িত, অর্থাৎ $\dom{f} = \Setabs{x \in A}{f(x) \fdefined}$।
\end{defn}

\begin{ex}
প্রত্যেক অপেক্ষক $f\colon A \to B$ একই সঙ্গে একটি আংশিক অপেক্ষক। যেসব আংশিক অপেক্ষক~$A$-এর সর্বত্র সংজ্ঞায়িত—অর্থাৎ এতক্ষণ আমরা যাদের শুধু অপেক্ষক বলেছি—তাদের \emph{সর্বত্র সংজ্ঞায়িত} অপেক্ষকও বলা হয়।
\end{ex}

\begin{ex}
$f(x) = 1/x$ দ্বারা নির্ধারিত আংশিক অপেক্ষক $f \colon \Real \pto \Real$, $x = 0$-তে অসংজ্ঞায়িত এবং অন্য সর্বত্র সংজ্ঞায়িত।
\end{ex}

\begin{prob}
$f\colon A \pto B$ দেওয়া থাকলে, নিচের নিয়মে আংশিক অপেক্ষক $g\colon B \pto
A$ সংজ্ঞায়িত করো: যেকোনো $y \in B$-এর ক্ষেত্রে, $f(x) = y$ হয় এমন একটি অনন্য $x \in A$ থাকলে $g(y) = x$; অন্যথায় $g(y) \fundefined$। দেখাও যে $f$ একৈক হলে সকল $x \in \dom{f}$-এর জন্য $g(f(x)) = x$ এবং সকল $y \in \ran{f}$-এর জন্য $f(g(y)) = y$ হয়।
\end{prob}

\begin{defn}[আংশিক অপেক্ষকের গ্রাফ]
ধরি, $f\colon A \pto B$ একটি আংশিক অপেক্ষক। নিচের সংজ্ঞা দ্বারা নির্ধারিত সম্পর্ক $R_f \subseteq A \times B$-কে~$f$-এর \emph{গ্রাফ} বলা হয়:
\[
R_f = \Setabs{\tuple{x,y}}{f(x) = y}.
\]
\end{defn}

\begin{prop}
ধরি, $R \subseteq A \times B$-এর এই ধর্ম আছে যে $Rxy$ ও $Rxy'$ হলেই $y = y'$ হয়। তাহলে $R$ হলো নিচের নিয়মে সংজ্ঞায়িত আংশিক অপেক্ষক $f\colon A \pto B$-এর গ্রাফ: $Rxy$ হয় এমন কোনো $y$ থাকলে $f(x) = y$; অন্যথায় $f(x) \fundefined$। যদি $R$ একই সঙ্গে \emph{সিরিয়াল} হয়, অর্থাৎ প্রত্যেক $x \in A$-এর জন্য এমন একটি $y \in B$ থাকে, যাতে $Rxy$ হয়, তবে $f$ সর্বত্র সংজ্ঞায়িত।
\end{prop}

\begin{proof}
ধরি, এমন একটি $y$ আছে, যাতে $Rxy$। যদি অন্য কোনো $y'
\neq y$-এর জন্য $Rxy'$ হতো, তবে $R$-এর উপর আরোপিত শর্ত লঙ্ঘিত হতো। সুতরাং $Rxy$ হয় এমন কোনো $y$ থাকলে, সেই $y$ অনন্য; তাই $f$ সুসংজ্ঞায়িত। স্পষ্টতই $R_f = R$ এবং~$R$ সিরিয়াল হলে $f$ সর্বত্র সংজ্ঞায়িত।
\end{proof}

\end{document}
